% ══════════════════════════════════════════════════════════════════════════════
%  EXPERIMENT 5 — SOLUBILITY AND TEMPERATURE                     (Lesson 13, Term 2 Week 5)
% ══════════════════════════════════════════════════════════════════════════════
\experimentdiv{Solubility and Temperature}{Making a saturated solution \textbullet{} Pooling class data \textbullet{} Drawing a solubility curve}{Lesson 13 \textbullet\ Term 2, Week 5 \textbullet\ Lab}\label{sec:5}
\lessoncodes{VC2S8U05, U06 \textbullet{} VC2S8I02, I04, I05, I06}

\vocabhead
\begin{ctab}
\begin{tabular}{|SW{4.2cm}|W{11.6cm}|}
  \hline\hd & \thead{Word} & \thead{Definition (copy from the slide)} \\ \hline
  & \scaf{\textbf{Saturated}} & \scaf{A solution that holds as much solute as it can at that temperature; any more stays undissolved.} \\ \hline
  & \textbf{Solubility} & \\ \hline
  & \textbf{Solubility curve} & \\ \hline
  & \textbf{Anomaly} & \\ \hline
  & \textbf{Line of best fit} & \\ \hline
\end{tabular}
\end{ctab}

\begin{lipanel}
\inquiryli{Measure precisely, pool class data and draw a line graph with a line of best fit.}
\understandli{Describe how the solubility of a solid changes with temperature.}
\sctitle
\begin{itemize}[leftmargin=1.4em, itemsep=0pt, topsep=1pt]
  \item I can tell when a solution is saturated.
  \item I can plot class data and draw a smooth line of best fit.
  \item I can spot an anomaly and suggest a cause.
\end{itemize}
\end{lipanel}

\begin{skillsfocus}
\focus{PC4}{Measure mass and temperature precisely.}
\focus{PC6}{Collect your group's result for the class.}
\focus{PA1}{Pool the class data in one table.}
\focus{PA2}{Plot a solubility curve.}
\focus{PA5}{Describe the trend.}
\focus{PA7}{Find and explain any anomalies.}
\focus{EV2}{Compare the class curve with a published one.}
\end{skillsfocus}

% ─── TECHNIQUE ────────────────────────────────────────────────────────────────
\exppart{Technique spotlight}{Blue}

\begin{technique}{Reaching saturation}{PC4}
Add the solid a little at a time and stir until each portion has dissolved. The solution is \textbf{saturated} when a portion will not dissolve, even after a full minute of stirring. Count only the portions that dissolved completely.
\begin{enumerate}[leftmargin=1.6em, itemsep=1pt, topsep=3pt]
  \item Weigh each portion on the balance to the nearest 0.1\,g.
  \item Keep the temperature steady in a water bath, and read the thermometer at eye level.
  \item Stir for the same time after every portion.
\end{enumerate}
\end{technique}

\begin{technique}{Plotting a line graph}{PA2}
\begin{enumerate}[leftmargin=1.6em, itemsep=1pt, topsep=3pt]
  \item Put the independent variable (temperature) on the horizontal axis and the dependent variable (mass dissolved) on the vertical axis. Give each a label and a unit.
  \item Choose an even scale that uses most of the grid: each square must stand for the same amount.
  \item Plot each point as a small neat cross ($\times$), and give the graph a title that names both variables.
  \item Finish it with one smooth \textbf{line of best fit} through the middle of the points, not a line joining each point to the next. Class data always has some scatter: joining the points would turn every small error into a kink, while a line of best fit shows the trend and lets you read values \emph{between} the points, as you will at 45\,\textdegree C.
\end{enumerate}
\end{technique}

\begin{workedex}{Choosing a scale}
The largest mass is 10.5\,g and the grid is 16 squares high. 16 squares $\times$ 1\,g would fit, but 2 squares per gram (0.5\,g a square) would not (that needs 21 squares). Use 1\,g per square.
\end{workedex}

\begin{wedo}{Choosing a scale}
The temperatures go from 20\,\textdegree C to 60\,\textdegree C and the grid is 20 squares wide. Choose a scale for the horizontal axis and explain your choice.
\answerbox{3cm}
\end{wedo}

% ─── PLAN AND SAFETY / METHOD ────────────────────────────────────────────────
\exppart{Plan, safety and method}{Teal}

\textbf{Aim:} To investigate the effect of changing the temperature of water on the mass of potassium nitrate that dissolves in it.

Each group tests one temperature, then the class pools its results.\quad Our temperature: \rule{2cm}{0.4pt}\,\textdegree C

\task Write a hypothesis: \emph{If the temperature of the water is increased, the mass of potassium nitrate that dissolves will \dots} Add a reason using the particle model. \skr{QP1}
\answerbox{4cm}

\subsubsection*{Equipment}
Boiling tube \textbullet{} 10\,mL measuring cylinder \textbullet{} electronic balance (0.1\,g) \textbullet{} weighing boats \textbullet{} spatula \textbullet{} stirring rod \textbullet{} thermometer \textbullet{} 250\,mL beaker as a water bath, with kettle water mixed to your temperature \textbullet{} potassium nitrate

\begin{center}
\begin{tikzpicture}[scale=0.8]
  \eqBeaker
  \fill[black!8] (0.02,0.02) rectangle (1.78,1.9); \draw[Ink, line width=0.8pt] (0.02,1.9) -- (1.78,1.9);
  \begin{scope}[shift={(0.6,0.35)}]\eqTestTube\end{scope}
  \fill[black!25] (0.62,0.65) arc[start angle=180,end angle=360,radius=0.28] -- (1.18,1.6) -- (0.62,1.6) -- cycle;
  \begin{scope}[shift={(1.35,0.5)}, scale=0.95]\eqThermometer\end{scope}
  \eqlab{(1.2,3.3)}{(3.0,3.6)}{boiling tube}
  \eqlab{(1.6,2.9)}{(3.0,2.9)}{thermometer}
  \eqlab{(0.9,1.2)}{(3.0,1.8)}{10\,mL water $+$ potassium nitrate}
  \eqlab{(1.78,0.6)}{(3.0,0.6)}{water bath at your temperature}
\end{tikzpicture}
\expfigcap{The boiling tube stands in a water bath to hold its temperature.}
\end{center}

\begin{infobox}{Safety}
Wear safety glasses. Potassium nitrate is an oxidiser: keep it away from flames, and wash your hands after the lesson. Hot water scalds: your teacher fills the water bath. Put leftover solutions in the labelled waste container, not down the sink.
\end{infobox}

\subsubsection*{Method}
\begin{enumerate}[itemsep=2pt]
  \item Measure 10\,mL of water into the boiling tube and stand it in the water bath. Wait until the thermometer reads your temperature.
  \item Weigh 0.5\,g of potassium nitrate and add it to the tube. Stir until it has all dissolved.
  \item Keep adding 0.5\,g portions, one at a time, stirring after each.
  \item Stop when a portion will not dissolve after a minute of stirring. Record the total mass that \emph{did} dissolve.
  \item Write your result on the class table on the board, and copy the whole class's results into Table 5.1.
\end{enumerate}

% ─── RESULTS ─────────────────────────────────────────────────────────────────
\exppart{Results}{Amber}

\subsubsection*{Table 5.1: Mass of potassium nitrate that dissolves in 10\,mL of water (class data) \skr{PA1}}
\begin{ctab}
\begin{tabular}{|SW{3.0cm}|C{1.9cm}|C{1.9cm}|C{1.9cm}|C{1.9cm}|C{2.6cm}|}
  \hline\hd & \thead{Temperature (\textdegree C)} & \thead{Group 1 (g)} & \thead{Group 2 (g)} & \thead{Group 3 (g)} & \thead{Group 4 (g)} & \thead{Average (g)} \\ \hline
  & \scaf{20} & \scaf{3.0} & \scaf{3.5} & \scaf{3.0} & \scaf{---} & \scaf{3.2} \\ \hline
  & 20 & & & & & \\ \hline
  & 30 & & & & & \\ \hline
  & 40 & & & & & \\ \hline
  & 50 & & & & & \\ \hline
  & 60 & & & & & \\ \hline
\end{tabular}
\end{ctab}

\needspace{20\baselineskip}
\task Plot the class averages as a line graph and draw a smooth line of best fit. This is a \textbf{solubility curve}. \skr{PA2}
\graphgrid{10cm}

% ─── ANALYSIS ────────────────────────────────────────────────────────────────
\exppart{Analysis and evaluation}{Coral}

\task Describe the trend shown by your solubility curve. \skr{PA5}
\answerbox{4cm}

\needspace{10\baselineskip}
\task Look at the results of each group in Table 5.1. Is there an anomaly? Circle it and suggest what might have caused it. \skr{PA7}
\answerbox{4cm}

\needspace{10\baselineskip}
\task Use your curve to estimate the mass of potassium nitrate that would dissolve in 10\,mL of water at 45\,\textdegree C. Show on the graph how you found it.
\answerbox{2cm}

\needspace{14\baselineskip}
\task The published solubility of potassium nitrate is shown below. Compare it with the class results. Are the class values higher or lower? Suggest a reason. \skr{EV2}
\begin{ctab}
\begin{tabular}{|L{5.4cm}|C{1.6cm}|C{1.6cm}|C{1.6cm}|C{1.6cm}|C{1.6cm}|}
  \hline\thead{Temperature (\textdegree C)} & 20 & 30 & 40 & 50 & 60 \\ \hline
  \thead{Mass in 10\,mL of water (g)} & 3.2 & 4.6 & 6.4 & 8.5 & 11.0 \\ \hline
\end{tabular}
\end{ctab}
\answerbox{4cm}

\needspace{10\baselineskip}
\task A saturated solution at 60\,\textdegree C is cooled to 20\,\textdegree C. Predict what you would see, and explain using your curve.
\answerbox{4cm}

\begin{yr8box}
\begin{keepext}
\exttask The solubility of table salt hardly changes between 0\,\textdegree C and 100\,\textdegree C. Sketch on your graph what its solubility curve would look like, and explain why cooling a salt solution is not a good way to get salt crystals.
\answerbox{5cm}
\end{keepext}
\end{yr8box}

% ─── SKILL CHECK ─────────────────────────────────────────────────────────────
\exppart{Skill check}{NavyBlue}

\begin{skillcheck}
\scrow{PC4}{I measured mass to 0.1\,g and kept the temperature steady.}
\scrow{PC6}{I collected a result the class could use.}
\scrow{PA1}{I pooled the class data and calculated averages.}
\scrow{PA2}{I plotted a line graph with a line of best fit.}
\scrow{PA5}{I can describe the trend in a solubility curve.}
\scrow{PA7}{I can find an anomaly and suggest a cause.}
\scrow{EV2}{I can compare class results with published values.}
\end{skillcheck}

\task Which skill will you focus on in Experiment 6? Explain what you will do differently.
\answerbox{5cm}
